
Considérons le problème à valeur initiale suivant :
\[
y' = 1 + \frac{y}{t}, \quad \text{pour } t \in [1,\ 2], \quad \text{avec } y(1) = 2.
\]

\begin{enumerate}
    \item Trouver la solution exacte $y = y(t)$ de ce problème.
    
    \item On souhaite approximer la solution en utilisant une méthode numérique fondée sur l'idée suivante :
    
    \medskip
    
    À partir d'un point $(t_i, y_i)$, on estime la valeur suivante $y_{i+1}$ par la formule :
    \[
    \begin{cases}
    t_{i+1} = t_i + h, \\
    y_{i+1} = y_i + h \cdot f(t_i, y_i),
    \end{cases}
    \]
    où $f(t, y) = 1 + \dfrac{y}{t}$ est la fonction définissant l'équation différentielle, et $h$ est le pas de discrétisation.
    
    \medskip
    
    En utilisant cette méthode avec un pas $h = 0{,}1$, approximer la valeur de la solution en $t = 1{,}2$.
\end{enumerate}